% =========================================================
% Figura: arquitectura de módulos del sistema de automatización
% Requiere en el preámbulo:
%   \usepackage{tikz}
%   \usetikzlibrary{positioning,arrows.meta,fit,backgrounds}
% =========================================================

\begin{figure}[htbp]
  \centering
  \begin{tikzpicture}[
      font=\footnotesize,
      modulo/.style={
        draw, rounded corners=2pt, align=center,
        text width=30mm, minimum height=11mm, inner sep=2pt},
      hw/.style={modulo, fill=black!6},
      capa/.style={
        draw=black!35, dashed, rounded corners=4pt, inner sep=4mm, inner ysep=5mm},
      etiqueta/.style={font=\scriptsize\itshape, text=black!55},
      flecha/.style={-{Stealth[length=2mm]}, thick},
      pasiva/.style={-{Stealth[length=2mm]}, densely dashed, black!55},
      node distance=8mm and 10mm
    ]

    % --- Presentación ---
    \node[modulo] (gui) {Interfaz gráfica\\\scriptsize eventos y presentación};

    % --- Coordinación ---
    \node[modulo, below=14mm of gui] (med)
      {Medición\\\scriptsize orquestación de la secuencia};
    \node[modulo, left=of med] (conn)
      {Conexión y monitoreo\\\scriptsize estado y reintentos};
    \node[modulo, right=of med] (trig)
      {Disparo\\\scriptsize por tiempo o temperatura};

    % --- Hardware ---
    \node[hw, below=14mm of med] (oscil)
      {Osciloscopio\\\scriptsize VXI-11 sobre Ethernet};
    \node[hw, left=of oscil] (laser)
      {Láser\\\scriptsize RS-232};
    \node[hw, right=of oscil] (temp)
      {Temperatura\\\scriptsize ESP32 y sensores};

    % --- Datos y seguridad ---
    \node[modulo, below=14mm of oscil] (store)
      {Almacenamiento\\\scriptsize registro de sesión};
    \node[modulo, left=of store] (safe)
      {Modo seguro\\\scriptsize solo actúa sobre el láser};
    \node[modulo, right=of store] (viz)
      {Visualización\\\scriptsize consulta y exportación};

    % --- Cajas de capa ---
    \begin{scope}[on background layer]
      \node[capa, fit=(conn)(med)(trig)] (cap_coord) {};
      \node[capa, fit=(laser)(oscil)(temp)] (cap_hw) {};
      \node[capa, fit=(safe)(store)(viz)] (cap_datos) {};
    \end{scope}

    \node[etiqueta, anchor=north west] at ([xshift=1mm,yshift=-0.5mm]cap_coord.north west) {coordinación};
    \node[etiqueta, anchor=north west] at ([xshift=1mm,yshift=-0.5mm]cap_hw.north west) {instrumentos};
    \node[etiqueta, anchor=north west] at ([xshift=1mm,yshift=-0.5mm]cap_datos.north west) {datos};

    % --- Flujo de control ---
    \draw[flecha] (gui) -- node[right, font=\scriptsize] {modo automático} (med);
    \draw[flecha] (gui.west) -- ++(-2.2,0) |- (conn.north);
    \draw[flecha] (med) -- (trig);
    \draw[flecha] (med) -- (oscil);
    \draw[flecha] (med.south west) -- (laser.north east);
    \draw[flecha] (med.south east) -- (temp.north west);
    \draw[flecha] (trig.south) -- (temp.north);
    \draw[flecha] (oscil) -- (store);
    \draw[flecha] (med.north west) -- ++(-0.35,0) -- ++(0,-0.4)
      -| ([xshift=-6mm]cap_hw.west) |- (safe.west);

    % --- Vigilancia pasiva ---
    \draw[pasiva] (conn.south) -- (laser.north);
    \draw[pasiva] (conn.south east) -- (oscil.north west);
    \draw[pasiva] (conn.east) -- ++(0.3,0) -- ++(0,-0.55)
      -| ([xshift=6mm]temp.north east) -- ([xshift=6mm]temp.north east|-temp.east) -- (temp.east);

    % --- Modo seguro ---
    \draw[flecha] (safe) -- (laser);

    % --- Consulta ---
    \draw[flecha] (store) -- (viz);

  \end{tikzpicture}
  \caption[Arquitectura de módulos del sistema de automatización]{%
    Arquitectura de módulos del sistema de automatización. Las flechas
    continuas indican control o transferencia de datos; las discontinuas,
    la vigilancia pasiva del estado de conexión. El módulo de medición
    coordina la secuencia y es el único que escribe en el registro de
    sesión. El modo seguro actúa exclusivamente sobre el láser: la
    configuración del osciloscopio no se modifica desde el programa.}
  \label{fig:arquitectura}
\end{figure}
